% ECONOMICS_PROBLEM: {"id": "O41-563", "title": "Demography and long-run growth", "jel_code": "O41", "source_id": "OP-0269"}
% !TeX program = xelatex
\documentclass[11pt,a4paper]{article}
\usepackage[margin=23mm]{geometry}
\usepackage{fontspec}
\setmainfont{Latin Modern Roman}
\usepackage{amsmath,amssymb,mathtools}
\usepackage{xcolor,enumitem,booktabs,longtable,array,xurl,hyperref}
\definecolor{muted}{HTML}{53636C}
\hypersetup{colorlinks=true,urlcolor=blue,linkcolor=blue}
\setlength{\parindent}{0pt}
\setlength{\parskip}{5pt}
\newcommand{\R}{\mathbb R}
\newcommand{\E}{\mathbb E}
\newcommand{\N}{\mathbb N}
\newcommand{\ind}{\mathbf 1}
\newcommand{\Var}{\operatorname{Var}}
\newcommand{\Cov}{\operatorname{Cov}}
\newcommand{\supp}{\operatorname{supp}}
\DeclareMathOperator*{\argmin}{arg\,min}
\DeclareMathOperator*{\argmax}{arg\,max}
\newcommand{\entrymeta}[3]{{\sffamily\small\color{muted}\textbf{\texttt{#1}}\quad|\quad #2\par #3\par}}
\newcommand{\Problemref}[1]{\texttt{#1}}
\newcommand{\JELref}[2]{\texttt{#2} (JEL \texttt{#1})}
\begin{document}
\section*{Demography and long-run growth}
\textbf{Problem ID:} \texttt{O41-563}\quad \textbf{JEL:} \texttt{O41}\par
% BEGIN ECONOMICS STATEMENT
\label{problem:OP-0269}
{\sffamily\small\color{muted}Original subject group: Growth and productivity\par}
\label{definitions:problem:30007550}
\textbf{Audit classification:} Background; proposed extension. The model implication is established; transporting elasticities to demographic contraction is editorial empirical work. Verification concerns the cited version, not first priority or proof that this exact editorial specification remains unsolved on October 8, 2026.
\subsection*{Definitions and precise research target}
\paragraph{Editorial mathematical specification.} Consider the semi-endogenous idea-production relation
\[
\dot A_t=\chi R_t^\lambda A_t^\phi,\qquad
R_t=N_t s_t h_t z_t,
\]
where \(A_t>0\) is useful knowledge, \(R_t\) effective research effort, \(\chi>0\), \(\lambda>0\), and \(\phi<1\). Population \(N_t\), researcher share \(s_t\), research skill \(h_t\), and technological augmentation \(z_t\) determine effort; output per person uses a stated production relation with \(A_t\). For fixed parameters and positive balanced effort growth \(g_R>0\), when a positive balanced-growth path exists, the model’s familiar implication is \(g_A=\lambda g_R/(1-\phi)\); establishing that algebraic relation is not the unresolved target.
Identify a credible joint set for \((\lambda,\phi)\) and future effort components under an explicitly declared declining-population scenario. Distinguish research quality, endogenous allocation into research, and AI augmentation from raw researcher headcounts. Use independent innovation measures and variation in research supply to assess whether historically estimated elasticities transport to demographic contraction. Compute corresponding bounds for long-run per-capita growth or for finite transition paths when balanced growth does not exist. State sensitivity to the production relation and to knowledge obsolescence. The target is empirical identification and scenario validity; demographics alone do not determine a model-free growth forecast, and known special-case demographic growth results must not be mislabeled unsolved.

For measurable locally bounded \(R_t\ge0\) and \(A_0>0\), the stated differential equation already yields \(A_t^{1-\phi}=A_0^{1-\phi}+(1-\phi)\chi\int_0^tR_u^\lambda\,du\). If \(R_t=R_0e^{g_Rt}\) with \(g_R<0\), this integral converges and knowledge approaches a finite positive level; one cannot extrapolate the positive balanced-growth formula to a negative constant knowledge-growth rate in this nonobsolescent model. If \(g_R=0\), knowledge rises polynomially and its proportional growth tends to zero. Declining population \(N\) can coexist with growing effective research effort if \(s,h,z\) rise sufficiently. A per-capita growth prediction additionally needs a production/accumulation equation, and elasticities of the research law require normalized knowledge measurement and excluded effort variation.
\subsection*{Variants and assumption changes}
The following are \emph{editorial variants}, not additional published open conjectures.
\begin{enumerate}
\item \emph{Nonobsolescent declining effort.} Hold researcher share, skill, and augmentation fixed and impose exponential population decline. Use the explicit solution to calculate finite transitions and the limiting knowledge stock; this conditional calculation is already answered by the model.
\item \emph{Obsolescence and augmentation.} Add \(-\delta_AA_t\), \(\delta_A>0\), to idea production and allow \(z_t\) to grow. Reidentify \((\lambda,\phi,\delta_A)\) and compare whether AI augmentation offsets shrinking research population; the previous integral solution no longer applies.
\end{enumerate}
\subsection*{References and source audit}
\begin{itemize}
\item Charles I. Jones. \emph{The Past and Future of Economic Growth: A Semi-Endogenous Perspective}. 2022. Locator: Publisher metadata and Abstract, final paragraph; no full article claim made. Publisher metadata/abstract checked; full article was not retrieved. \url{https://www.annualreviews.org/content/journals/10.1146/annurev-economics-080521-012458}.
\end{itemize}
The model implication is established; transporting elasticities to demographic contraction is editorial empirical work.
\begin{enumerate}\raggedright
\item\hypertarget{definitions:source:30007550:1}{} Charles I. Jones. 2022. \emph{The Past and Future of Economic Growth: A Semi-Endogenous Perspective}. Publisher metadata and Abstract, final paragraph; no full article claim made.
\url{https://www.annualreviews.org/content/journals/10.1146/annurev-economics-080521-012458}.
DOI: \url{https://doi.org/10.1146/annurev-economics-080521-012458}.
\end{enumerate}

\subsection*{Common conventions: Source mathematical conventions}

These conventions apply to each entry unless explicitly replaced.

\textbf{Models and identification.} A primitive model \(m\) specifies all probability laws, feasible choices, information, equilibrium and measurement rules used by its target. The admissible class \(\mathcal M\) is fixed independently of outcomes. If \(P_O(m)\) is its observable probability law and \(\tau(m)\) the target, its identified set at observable law \(P\) is
\[
 \mathcal I_\tau(P)=\{\tau(m):m\in\mathcal M,\ P_O(m)=P\}.
\]
Point identification means that this set is a singleton. An empty set signals incompatibility of data and assumptions. Coordinate bounds need not describe the joint set. Bounds are sharp only if every retained target is attainable or arbitrarily approachable by compatible models. Finite bounds require explicit restrictions on unobserved quantities; unrestricted confounding, hidden wealth, extrapolation or arbitrary equilibrium selection can yield uninformative sets.

\textbf{Causal interventions.} A potential outcome \(Y_i(p)\) is the outcome under a complete policy regime \(p\), including timing, financing, information and market spillovers. The target population and units are fixed before treatment. With interference, use the complete assignment vector or a justified exposure mapping; individual no-interference notation alone cannot identify an economy-wide intervention. A difference of mean potential outcomes does not require the cross-world joint distribution. Conditioning on post-treatment migration, survival, enrollment or employment changes the estimand and needs separate justification.

\textbf{Statistical statements.} An estimator, confidence set, forecast or test specifies a sampling law, model class and loss/coverage criterion. Uniform \((1-\alpha)\)-coverage means \(\inf_{m\in\mathcal M}\Pr_m\{\tau(m)\in C_n\}\geq1-\alpha\), or an explicitly asymptotic version. The whole parameter space trivially has coverage, so informativeness requires a stated width, risk or power objective. A forecasting improvement is relative to a named benchmark, information set and evaluation population.

\textbf{Equilibrium and optimization.} An equilibrium comprises feasible plans, optimality, beliefs where needed, budget constraints and all market/resource clearing. The primitives or a stated benchmark define each correspondence \(\mathcal E(p;m)\). Nonemptiness is assumed or proved, never inferred just from compact policy sets. A selected-equilibrium target uses a declared selection rule; a robust target quantifies over all permitted equilibria. Use suprema or infima unless attainment is established. An \(\varepsilon\)-optimal policy is feasible and within \(\varepsilon\) of the finite optimum value. Report an empty feasible set separately.

\textbf{Welfare and accounting.} Normative utility, interpersonal normalization, social weights, numeraire, discounting and the use of public receipts are primitives. Behavioral utility need not coincide with the welfare criterion. Household welfare includes ownership income and rents once; adding profits again to compensating variation generally double-counts them. A monetary equivalent is defined by an explicit transfer and price convention, with monotonicity and range conditions. Average effects, mechanism contributions, transfers and welfare losses are different targets. Infinite sums require convergence and dynamic feasible plans require terminal or no-Ponzi conditions.

\textbf{Source versions.} A working-paper date, revision date and journal publication year are distinguished. A DOI for a journal version is not automatically the DOI of an earlier manuscript. Locators refer to the stated version and printed pages where specified. A metadata-only check does not verify a claim about a full-text conclusion. Every entry's source subsection narrows the provenance claim accordingly.
% END ECONOMICS STATEMENT
\end{document}
